%% pernyataan.tex
%% Lembar Pernyataan Orisinalitas Tesis
%% Font: Times New Roman 12pt
%%
%% EDIT DATA DI BAWAH INI SESUAI DATA TESIS ANDA

\newpage
\thispagestyle{plain}

\begin{center}
    {\fontsize{14}{16}\selectfont\bfseries {PERNYATAAN ORISINALITAS}}
\end{center}

\vspace{0.2cm}

\noindent
Yang bertanda tangan di bawah ini,

\vspace{0.5cm}

%% ===== DATA MAHASISWA DAN PEMBIMBING =====
\noindent
\begin{tabular}{@{}p{4.9cm}@{}p{0.5cm}@{}p{9cm}@{}}
    Nama Mahasiswa (NRP) & : & \\
    Dosen Pembimbing 1 (NIP) & : & \\
    Dosen Pembimbing 2 (NIP) & : & \\
    Program Studi & : & S2 Teknik Informatika \\
    Departemen & : & Departemen Teknik Informatika \\
    Fakultas & : & Fakultas Teknologi Elektro dan Informatika Cerdas \\
\end{tabular}

\vspace{0.5cm}

%% ===== PERNYATAAN =====
\noindent
Dengan ini menyatakan bahwa Tesis yang berjudul ``Judul Tesis'' adalah hasil karya sendiri, bersifat orisinal, dan ditulis dengan mengikuti kaidah penulisan ilmiah.

\vspace{0.3cm}

\noindent
Apabila di kemudian hari ditemukan ketidaksesuaian dengan pernyataan ini, maka saya bersedia menerima sanksi sesuai dengan ketentuan yang berlaku di Institut Teknologi Sepuluh Nopember.

\vspace{1.5cm}

%% ===== TANDA TANGAN MAHASISWA =====
\begin{flushright}
\begin{tabular}{@{}c@{}}
    Surabaya, dd mm \underline{yyyy} \\
    Mahasiswa, \\[1.5cm]
    TTD \underline{Bermaterai} \\[0.5cm]
    Nama Mahasiswa \\
    NRP.
\end{tabular}
\end{flushright}

\vspace{0.3cm}

%% ===== MENGETAHUI =====
\begin{center}
Mengetahui,
\end{center}

\vspace{0.3cm}

%% ===== TANDA TANGAN DOSEN PEMBIMBING =====
\noindent
\begin{tabular}{@{}p{7cm}@{}p{7cm}@{}}
    \centering Dosen Pembimbing 1, & \centering Dosen Pembimbing 2, \tabularnewline
    & \\[1.5cm]
    \centering Nama & \centering Nama \tabularnewline
    \centering NIP. & \centering NIP. \tabularnewline
\end{tabular}